\documentclass[10pt,a4paper]{amsart}
\usepackage[margin=19mm]{geometry}
\usepackage{amsmath,amssymb,mathtools}
\usepackage[scr=rsfs,cal=euler]{mathalfa}
\usepackage{xfrac}
\usepackage[unicode,hidelinks,bookmarks=false]{hyperref}
\newcommand{\ie}{i.e.\ }
\newcommand{\cf}{cf.\ }
\newcommand{\resp}{respectively\ }
\newcommand{\Subsection}{\subsection}
\providecommand{\nobreakdash}{\nobreak\hbox{-}}
\setlength{\emergencystretch}{2em}
\multlinegap0pt
\usepackage{amssymb}
\newtheorem{lem}{Lemma}
\newcommand{\I}{\mathcal{I}}
\newcommand{\K}{\mathbb{K}}
\newcommand{\PP}{\mathcal{P}}
\newcommand{\Z}{\mathcal{Z}}
\title{Retract rational varieties are uniformly retract rational}
\author{Juliusz Banecki}
\date{Source: arXiv:2411.17892v4 (CC BY 4.0)}
\begin{document}
\maketitle

\noindent\emph{Source and licence.} Retract rational varieties are uniformly retract rational. Version arXiv:2411.17892v4 (CC BY 4.0), distributed by arXiv under the Creative Commons Attribution 4.0 International licence, http://creativecommons.org/licenses/by/4.0/, as recorded in the arXiv licence metadata for this exact version and shown on the abstract page https://arxiv.org/abs/2411.17892v4. Full licence notice: \texttt{licenses/arxiv-2411.17892-CC-BY-4.0.txt}.\par\smallskip

\noindent\textit{Editorial scope.} Counted unit: the article's lem generic (source0.tex lines 273--293). The article's complete original proof of it is delivered with it (source0.tex lines 225--269, one \texttt{proof} environment of 45 lines, which the article places away from the statement). No mathematics is written, repaired or re-derived: the statement and the proof are the article's own lines, in the article's own notation, and no retained line is substituted or re-flowed.  No further source-local context is needed: every cross-reference of every retained line resolves inside the two delivered spans.  Nothing was omitted from any retained span, so the package carries no omission record.  No retained line contains a citation, so the wrapper carries no bibliography and no bibliography entry.  No retained line contains a figure environment, an image-inclusion command or drawing code, so no image is delivered and the package declares no asset dependency.  Editorial formatting only, in addition: an AMS \texttt{amsart} preamble that restates the article's own declarations and macros used by the retained lines, namely the environments \texttt{lem} on separate counters, together with the standard packages those lines need.  The retained environments print in this document as Lemma 1 in order of appearance, whereas the article prints the same items with the numbering of its own full document.  The complete native source of the article as distributed in the arXiv e-print for this exact version is bundled unchanged as \texttt{sources/169188/source0.tex}.
\begin{lem}\label{generic}
Let $X\subset \K^n$ be a Zariski closed irreducible subvariety of $\K^n$ of dimension $m$ containing the origin and let $I\subset \PP(X)$ be a nonzero ideal with $0\in\Z(I)$. Then, a generic element $W$ of the Grassmannian $\mathrm{Gr}_{n-m}(\K^n)$ satisfies
\begin{equation*}
   W^{\overline{\K}}\cap \Z^{\overline{\K}}(I)=\{0\}.  
\end{equation*}
\begin{proof}
It suffices to consider the case when $\K$ itself is algebraically closed, since $\mathrm{Gr}_{n-m}(\K^n)$ is irreducible and it is Zariski dense in $\mathrm{Gr}_{n-m}(\overline{\K}^n)$. 

 

Let $Z:=\Z(I)$ and consider the variety
\begin{equation*}
    Y:=\{(V,x)\in\mathrm{Gr}_{n-m}(\K^n)\times (Z\backslash \{0\}):x\in V\}. 
\end{equation*}
Each fiber under the natural projection $\pi_2:Y\rightarrow Z\backslash \{0\}$ is isomorphic to the lower dimensional Grassmannian $\mathrm{Gr}_{n-m-1}(\K^{n-1})$. Since $I$ is nonzero it folows that $\dim Z\leq m-1$, so 
\begin{equation*}
    \dim Y\leq \dim \mathrm{Gr}_{n-m-1}(\K^{n-1})+\dim Z\leq(n-m)m-1<\dim \mathrm{Gr}_{n-m}(\K^n).
\end{equation*}
This shows that a generic element of $\mathrm{Gr}_{n-m}(\K^n)$ does not lie in the image of the projection $\pi_1:Y\rightarrow \mathrm{Gr}_{n-m}(\K^n)$, which finishes the proof.
\end{proof}
\end{lem}

\begin{proof}
Without loss of generality we can assume $x_0=0\in\K^n$. Let $W$ be a generic vector subspace of $\K^n$ of codimension $m$. 


We can assume that it intersects $X$ transversally at zero. Moreover, by a simple dimension counting argument (see Lemma \ref{generic} below), we can assume that $\Z^{\overline{\K}}(I)\cap W^{\overline{\K}}=\{0\}$. 

Let $x_1,\dots,x_n$ be a basis of the dual space of $\K^n$ such that $W$ is given by the equation $x_{n-m+1}=\dots=x_n=0$. From Noether normalisation lemma we can assume that
\begin{equation}\label{finite_ext_eq}
	\PP(X) \text{ is integral over }\K[\bar{x}_1,\dots,\bar{x}_m],
\end{equation} 
where $\bar{x}_i$ denotes the restriction $x_i+\I(X)\in\PP(X)$ of $x_i$ to $X$.

For $1\leq i\leq k$ let $t_i:=x_i+\I(W)\in\PP(W)$. The linear system of parameters $t_1,\dots,t_{n-m}$ allows for a natural identification of $W$ with $\K^{m-n}$. Under this identification define
\begin{align*}
	&\sigma:X\times W\rightarrow \K^n,\\
	&\sigma(x,v):=x+v.
\end{align*}  
It is clear that under such a definition of $\sigma$, the properties $(1)$ and $(2)$ are satisfied. 


Consider the ring extension $\PP(\K^n)\subset \PP(X\times W)$ induced by $\sigma$. Identifying $\PP(X)$ with its image under the natural inclusion $\PP(X)\subset \PP(X\times W)$ for $n-m< i \leq n$ we have that $\bar{x}_i= x_i\in \PP(\K^n)$, so from \eqref{finite_ext_eq} we deduce that all the elements of $\PP(X)$ are integral over $\PP(\K^n)$. Since for $1\leq i \leq m$ we have $t_i+\bar{x}_i= x_i\in \PP(\K^n)$, it follows that the entire ring $\PP(X\times W)$ is integral over $\PP(\K^n)$. As it is finitely generated as an algebra, it is finitely generated as a module over $\PP(\K^n)$. 

Denote by $\mathfrak{n}$ the maximal ideal of functions vanishing at the point $(0,0)$ in $\PP(X\times W)$ and by $\mathfrak{m}=\mathfrak{n}\cap \PP(\K^n)$ the maximal ideal of functions vanishing at the point $0$ in $\PP(\K^n)$. 

Note that if $(x,v)\in \Z^{\overline{\K}}(I(t_1,\dots,t_{n-m}))\subset X^{\overline{\K}}\times W^{\overline{\K}}$ is a point with $\sigma^{\overline{\K}}(x,v)=0$, then either $x\in \Z^{\overline{\K}}(I)\cap W^{\overline{\K}}=\{0\}$, so $x=0$ and $v=0$, or $v\in \Z^{\overline{\K}}(t_1,\dots,t_m)=\{0\}$, so again $v=0$ and $x=0$. This shows that 
\begin{equation*}
    \Z^{\overline{\K}}(\mathfrak{m}\PP(X\times W) +I(t_1,\dots,t_{n-m}))=\{(0,0)\}.
\end{equation*}
In algebraic terms this means that $\mathfrak{n}$ is the only associated prime of $\mathfrak{m}\PP(X\times W) +I(t_1,\dots,t_{n-m})$.

Now, from point $(2)$ we deduce that elements of $\mathfrak{m}$ generate the maximal ideal of the local ring $\PP(X\times W)_{\mathfrak{n}}$. This means that for every $f\in\mathfrak{n}$ there is $q\in\PP(X\times W)\backslash\mathfrak{n}$ such that
\begin{equation*}
	qf\in \mathfrak{m}\PP(X\times W)\subset \mathfrak{m}\PP(X\times W) +I(t_1,\dots,t_{n-m}).
\end{equation*}
As the latter ideal is $\mathfrak{n}$-primary we get that 
\begin{equation}\label{ideal_equal_eq}
	\mathfrak{m}\PP(X\times W) +I(t_1,\dots,t_{n-m})=\mathfrak{n}.
\end{equation}

The ideal $I(t_1,\dots,t_{n-m})$ is a submodule of $\PP(X\times W)$. Consider the quotient module $M:=\PP(X\times W)/(I(t_1,\dots,t_{n-m})+\PP(\K^n))$. From \eqref{ideal_equal_eq} we deduce that 
\begin{equation*}
	M\subset M\mathfrak{m}.
\end{equation*}
Hence, it follows from Nakayama's lemma that $M_{\mathfrak{m}}=0$. This is equivalent to $(3)$.
\end{proof}

\end{document}
